
\section{典范模型、形式超概形与仓西映射}
\label{sec:5-4}

\subsection{典范模型、仓西映射与定界余链}
\label{subsec:5-4-1}

本节将完成导论中定理 A 和 F 的证明。
也就是说，我们要证明：
先前在上链复形上构造的滤过 $A_\infty$ 代数结构
与滤过 $A_\infty$ 双模结构，
会在上同调群上诱导滤过 $A_\infty$ 结构。
我们还把定界余链的集合
$\widehat{\mathcal M}(C)$
描述为一个形式映射，即仓西映射的零点集；
这也就是导论中的定理 M。
定理 N 以及推论 O、P
则将在第 5.4.6 小节证明。
在整个第 5.4 节中，
我们假设滤过 $A_\infty$ 代数
及滤过 $A_\infty$ 双模都是有隙的，
并假设 $R$ 是域；
不过，只有在证明引理 5.4.28
和 5.4.51 时才会用到 $R$ 为域这一条件。

先从下述定理 5.4.1 出发。
它并非新结果，已在
\textbf{[Kad82]} 中证明
（亦见 \textbf{[Mer99, KoS01]}；
其中以弱同伦等价代替同伦等价）。
这些文献所用的符号约定与本书不同；
另见 \textbf{[Fuk03II]}。

\begin{theorem}{5.4.1}\label{thm:5-4-1}
任意无滤过 $A_\infty$ 代数
$(\overline C,\overline{\mathfrak m})$
都同伦等价于一个
$A_\infty$ 代数
$(\overline C',\overline{\mathfrak m}')$，
且 $\overline{\mathfrak m}'_1=0$。
\end{theorem}

这里不证明定理 5.4.1，
而将在第 5.4.4 小节证明
它对于滤过 $A_\infty$ 代数的对应版本，
即定理 5.4.2。
该版本此前尚未见于文献。

\begin{theorem}{5.4.2}\label{thm:5-4-2}
任意有隙滤过 $A_\infty$ 代数
$(C,\mathfrak m)$
都同伦等价于一个有隙滤过
$A_\infty$ 代数
$(C',\mathfrak m')$，
其 $R$-约化
$(\overline C',\overline{\mathfrak m}')$
满足
$\overline{\mathfrak m}'_1=0$。
而且，此同伦等价可取为一个有隙
$A_\infty$ 同态。
\end{theorem}

定理 W 立即由定理 5.4.2 推出。

\newpage
\liftIndex{I208}\liftIndex{I209}
\begin{definition}{5.4.3}\label{def:5-4-3}
\begin{enumerate}
\item
若一个无滤过 $A_\infty$ 代数满足
$\overline{\mathfrak m}_1=0$，
就称它为\emph{典范的}。
若一个滤过 $A_\infty$ 代数满足
$\overline{\mathfrak m}_1=0$，
也称它为\emph{典范的}。
滤过 $A_\infty$ 代数 $C$ 的
\emph{典范模型}，
定义为与 $C$ 同伦等价的
典范滤过 $A_\infty$ 代数。

\item\liftIndex{I206}
设 $(C,\mathfrak m)$、
$(C',\mathfrak m')$
为滤过 $A_\infty$ 代数，
$\widehat d,\widehat d'$
分别是由其结构在
$\widehat BC[1]$、
$\widehat BC'[1]$
上诱导的余导子。
若存在滤过 $A_\infty$ 同态
$\mathfrak f\colon C\to C'$，
使诱导的余代数同态
$\widehat{\mathfrak f}\colon
\widehat BC[1]\to\widehat BC'[1]$
为双射，且
$\widehat{\mathfrak f}\circ\widehat d
=\widehat d'\circ\widehat{\mathfrak f}$、
其逆映射满足
$\widehat d\circ\widehat{\mathfrak f}^{-1}
=\widehat{\mathfrak f}^{-1}\circ\widehat d'$，
便称 $(C,\mathfrak m)$
与 $(C',\mathfrak m')$
\emph{同构}。
此时称 $\mathfrak f$
为滤过 $A_\infty$ 代数的一个同构。
无滤过版本也可同样定义。

\item\liftIndex{I207}
设 $(D,\mathfrak n)$、
$(D',\mathfrak n')$
为 $(C_1,C_0)$-滤过
$A_\infty$ 双模，
$\widehat d,\widehat d'$
为其诱导的双余导子。
若存在覆盖恒等同态的
滤过 $A_\infty$ 双模同态
$\varphi\colon D\to D'$，
使诱导的余模同态
\[
\widehat\varphi\colon
\widehat BC_1[1]\widehat\otimes
D[1]\widehat\otimes
\widehat BC_0[1]
\longrightarrow
\widehat BC_1[1]\widehat\otimes
D'[1]\widehat\otimes
\widehat BC_0[1]
\]
为双射，并且其逆映射满足
$\widehat d\circ\widehat\varphi^{-1}
=\widehat\varphi^{-1}\circ\widehat d'$，
便称 $(D,\mathfrak n)$
与 $(D',\mathfrak n')$
\emph{同构}。
此时称 $\varphi$
为滤过 $A_\infty$ 双模的一个同构。
无滤过版本也可同样定义。
\end{enumerate}
\end{definition}

\begin{remark}{5.4.4}\label{rem:5-4-4}
孔采维奇在 $L_\infty$ 代数的语境中，
把一个相近的概念称为\emph{极小的}
\textbf{[Kon03]}。
但在交换 D.G.A. 理论中，
沙利文用“极小”指称另一个概念
\textbf{[Sul78]}。
事实上，若把德拉姆复形
$(\Omega(M),d)$
看成 $A_\infty$ 代数，
它的典范模型作为向量空间
就是上同调群 $H(M;\mathbb R)$；
而沙利文意义下的极小模型\liftIndex{I342}，
除非 $M$ 是形式的，
否则不同于 $H(M;\mathbb R)$。
为避免混淆，这里采用“典范”一词。
注意，若 $C$ 是典范 $A_\infty$ 代数，
其 bar 复形 $BC[1]$
的对偶 $BC[1]^*$
是沙利文意义下的极小 D.G.A.。

还需指出，我们证明滤过 $A_\infty$
代数各项代数结果的顺序，
与 \textbf{[MSS02, Lef03]}
等文献大不相同。
那些文献先证明定理 5.4.1 的对应结果，
再用它来定义同伦等价等概念。
作者认为，这受到有理同伦论
及 D.G.A. 同伦论传统的影响。
在 D.G.A. 范畴中，
定理 4.2.45、命题 4.2.37
等若干基本结果只对极小 D.G.A. 成立，
因而必须在理论开始时
先建立极小模型的存在性。
而在滤过 $A_\infty$ 代数的范畴中，
这些结果对\emph{任意}
$A_\infty$ 代数都成立。
因此，典范滤过 $A_\infty$ 代数
在这一理论中，并不扮演
极小 D.G.A. 在 D.G.A. 同伦论中
所扮演的相同角色。
\end{remark}

下面的命题说明“典范”这一名称的根据。

\begin{proposition}{5.4.5}\label{prop:5-4-5}
两个典范的滤过或无滤过
$A_\infty$ 代数之间的
任意同伦等价，都是同构。
\end{proposition}

\begin{proof}
设 $\mathfrak f\colon C_1\to C_2$
为两个典范滤过 $A_\infty$
代数之间的同伦等价。
它在约化上诱导同伦等价
$\overline{\mathfrak f}\colon
\overline C_1\to\overline C_2$。
各 $\overline C_i$ 仍是典范的，
因而同构于自身的
$\overline{\mathfrak m}_1$-上同调。
所以
$\overline{\mathfrak f}_1\colon
\overline C_1\to\overline C_2$
是 $R$-模同构。
利用数目滤过，可按前文的标准论证
证明
$\widehat{\overline{\mathfrak f}}\colon
B\overline C_1[1]\to
B\overline C_2[1]$
是 $R$-模同构。
再利用能量滤过，
即可求得
$\widehat{\mathfrak f}\colon
BC_1[1]\to BC_2[1]$
的逆。
命题得证。
\end{proof}

由此可把定理 5.4.2 重述如下。

\begin{theorem}{5.4.2$'$}\label{thm:5-4-2-prime}
任意有隙滤过 $A_\infty$ 代数
$(C,\mathfrak m)$
都有一个典范模型，
并且它在同构意义下唯一。
同伦等价可取为有隙
$A_\infty$ 同态。
若 $(C,\mathfrak m)$ 同伦含幺，
则其典范模型含幺。
\end{theorem}

定理 5.4.2$'$
中关于幺元的断言，
并不直接由定理 5.4.2 得出。
我们将在第 5.4.4 小节
证明定理 5.4.2 后立即证明这一断言。
把以上结果用于几何情形，
立即得到下述推论。

\begin{corollary}{5.4.6}\label{cor:5-4-6}
对辛流形中每个相对 spin 的
拉格朗日子流形 $L$，
可关联一个含幺的典范滤过
$A_\infty$ 代数
$(C_{\mathrm{can}}(L;\Lambda_{0,\mathrm{nov}}),
\mathfrak m_{\mathrm{can}})$。
在含幺同构意义下，
它不依赖于殆复结构及构造中
所作的任何其他选择。
\end{corollary}

注意
$C_{\mathrm{can}}(L;\Lambda_{0,\mathrm{nov}})
\cong H(L;\Lambda_{0,\mathrm{nov}})$。
因此，推论 5.4.6 立即蕴涵定理 A。

下面利用典范模型定义仓西映射，
并研究它同定界余链模空间
$\widehat{\mathcal M}(C)$ 的关系。

从现在起假设
$H_*(\overline C,\overline{\mathfrak m}_1)$
是有限生成 $R$-模。
若 $\overline C$ 为典范的，
则 $\overline C$ 本身有限生成，
从而 $C$ 是有限生成
$\Lambda_{0,\mathrm{nov}}$-模。
这里用到我们一贯的假设：
$C$ 是自由
$\Lambda_{0,\mathrm{nov}}$-模，
且
$C=\overline C\otimes_R
\Lambda_{0,\mathrm{nov}}$。

\begin{definition}{5.4.7}\label{def:5-4-7}
设 $(C,\mathfrak m)$ 为滤过
$A_\infty$ 代数。
取 $m$ 个偶次数变量
$s_i$（$i=1,\ldots,m$）
和 $\ell$ 个奇次数变量
$h_i$（$i=1,\ldots,\ell$）。
\emph{超形式幂级数环}
$\Lambda_{0,\mathrm{nov}}
[[s_1,\ldots,s_m;h_1,\ldots,h_\ell]]$
是下述形式幂级数所成的环：
\[
\Lambda_{0,\mathrm{nov}}
[[s_1,\ldots,s_m;h_1,\ldots,h_\ell]]
=
\left\{
\sum_{k_1=0}^{\infty}\cdots
\sum_{k_m=0}^{\infty}
a_{k_1,\ldots,k_m}
s_1^{k_1}\cdots s_m^{k_m}
\right\},
\]
其中
$a_{k_1,\ldots,k_m}
\in\Lambda_{0,\mathrm{nov}}
\otimes_R\Omega(h_1,\ldots,h_\ell)$。
$\Omega(h_1,\ldots,h_\ell)$
表示由 $h_1,\ldots,h_\ell$
生成的 $R$ 上外代数。
环 $R[[s_1,\ldots,s_m;
h_1,\ldots,h_\ell]]$
也以同样方式定义。
\end{definition}

现设 $(C,\mathfrak m)$
为典范滤过 $A_\infty$ 代数。
取 $C^{\mathrm{odd}}
=C[1]^{\mathrm{even}}$
的一组基
$\mathbf v_1,\ldots,\mathbf v_m$，
以及 $C^{\mathrm{even}}$
的一组基
$\mathbf w_1,\ldots,\mathbf w_\ell$。
对 $1\le i\le m+\ell$，
令
$\pi_i\colon C\to\Lambda_{0,\mathrm{nov}}$
为坐标投影：
$i\le m$ 时取
$\mathbf v_i$ 的坐标，
$i>m$ 时取
$\mathbf w_{i-m}$ 的坐标。
若 $C$ 含幺，
约定 $\mathbf w_\ell=\mathbf e$，
其中 $\mathbf e$ 是幺元。

\begin{definition}{5.4.8}\label{def:5-4-8}
令
\[
z=s_1\mathbf v_1+\cdots+s_m\mathbf v_m
+h_1\mathbf w_1+\cdots+h_\ell\mathbf w_\ell.
\]
对 $i=1,\ldots,m+\ell$，定义
\begin{equation}
\begin{aligned}
P_i(s_1,\ldots,s_m;
h_1,\ldots,h_\ell)
&=(\pi_i\circ\mathfrak m)(\exp z)\\
&\in
\Lambda_{0,\mathrm{nov}}
[[s_1,\ldots,s_m;h_1,\ldots,h_\ell]].
\end{aligned}
\tag{5.4.9}\label{eq:5-4-9}
\end{equation}
\end{definition}

下面解释 (5.4.9) 右边的意义。
记
$\mathbf u_i^0=\mathbf v_i$、
$\mathbf u_i^1=\mathbf w_i$，
$u_i^0=s_i$、
$u_i^1=h_i$。
把指数函数形式展开后，
得到以下形式幂级数。
为使每个符号都有定义，
记 $\mathcal I_k$ 为
满足 $\epsilon_j\in\{0,1\}$，
且在 $\epsilon_j=0$ 时
$1\le a_j\le m$、
在 $\epsilon_j=1$ 时
$1\le a_j\le\ell$ 的
$k$ 元指标组 $(a_j,\epsilon_j)_{j=1}^k$。
\translatornote{原书实体页 342
在式 (5.4.10) 中统一令
$a_j\le\max\{m,\ell\}$；
但两类变量的定义域分别
只有 $m$ 项与 $\ell$ 项。
此处逐奇偶类型限制指标，
见疑误 A-0314。}
\begin{equation}
\begin{aligned}
\sum_{k\ge0}
\ \sum_{(a,\epsilon)\in\mathcal I_k}
&(-1)^{
\sum_{i<j}\epsilon_i\epsilon_j
+\sum_i\epsilon_i}
u_{a_1}^{\epsilon_1}
\cdots
u_{a_k}^{\epsilon_k}\\
&{}\cdot
(\pi_i\circ\mathfrak m)
\bigl(
\mathbf u_{a_1}^{\epsilon_1},
\ldots,
\mathbf u_{a_k}^{\epsilon_k}
\bigr).
\end{aligned}
\tag{5.4.10}\label{eq:5-4-10}
\end{equation}

\begin{definition}{5.4.11}\label{def:5-4-11}
把
$P_i(s_1,\ldots,s_m;
h_1,\ldots,h_\ell)$
称为典范滤过 $A_\infty$
代数 $(C,\mathfrak m)$
的\emph{超仓西映射}。
令所有 $h_i=0$，
得到
$P_i^{\mathrm{bos}}
(s_1,\ldots,s_m)$；
称它为 $(C,\mathfrak m)$
的\emph{仓西映射}。
\liftIndex{I210}
\translatornote{原书实体页 342
在本定义中把
$P_i$ 的偶变量写到 $s_\ell$、
奇变量写到 $h_m$，
纯偶限制也写到 $s_\ell$；
本页定义 5.4.7 与式 (5.4.9)
均固定应为 $s_m,h_\ell$。
见疑误 A-0315。}
定义
\[
\begin{aligned}
\mathfrak M(C,\mathfrak m)
&=
\Lambda_{0,\mathrm{nov}}
[[s_1,\ldots,s_m;h_1,\ldots,h_\ell]]
/(P_1,\ldots,P_{m+\ell}),\\
\mathfrak M(C,\mathfrak m)^{\mathrm{bos}}
&=
\Lambda_{0,\mathrm{nov}}
[[s_1,\ldots,s_m]]
/(P_{m+1}^{\mathrm{bos}},
\ldots,P_{m+\ell}^{\mathrm{bos}}).
\end{aligned}
\]
这里
$(P_1,\ldots,P_{m+\ell})$
指由这些元素生成的理想。
若 $C$ 含幺，则定义
\[
\begin{aligned}
\mathfrak M_{\mathrm{weak}}(C,\mathfrak m)
&=
\Lambda_{0,\mathrm{nov}}
[[s_1,\ldots,s_m;h_1,\ldots,h_\ell]]
/(P_1,\ldots,P_{m+\ell-1}),\\
\mathfrak M_{\mathrm{weak}}(C,\mathfrak m)^{\mathrm{bos}}
&=
\Lambda_{0,\mathrm{nov}}
[[s_1,\ldots,s_m]]
/(P_{m+1}^{\mathrm{bos}},
\ldots,P_{m+\ell-1}^{\mathrm{bos}}).
\end{aligned}
\]
注意
$P_{m+\ell}^{\mathrm{bos}}
=P_{m+\ell}$
取出幺元 $\mathbf e$ 的系数。
\end{definition}

对未必典范的一般滤过
$A_\infty$ 代数 $(C,\mathfrak m)$，
定义它的（超）仓西映射为
其典范模型的（超）仓西映射。
典范模型作为
$\Lambda_{0,\mathrm{nov}}$-模，
同构于 $(C,\overline{\mathfrak m}_1)$
的上同调，
故仓西映射可视为
定义在该上同调上的一个形式映射\liftIndex{I018}。
还要指出，
$\mathfrak M(C,\mathfrak m)^{\mathrm{bos}}$
是交换
$\Lambda_{0,\mathrm{nov}}$-代数，
而 $\mathfrak M(C,\mathfrak m)$
是超交换
$\Lambda_{0,\mathrm{nov}}$-代数。

\begin{theorem}{5.4.12}\label{thm:5-4-12}
把 $\mathfrak M(C,\mathfrak m)$
看成超交换
$\Lambda_{0,\mathrm{nov}}$-代数，
把
$\mathfrak M(C,\mathfrak m)^{\mathrm{bos}}$
看成交换
$\Lambda_{0,\mathrm{nov}}$-代数；
则二者的同构类
分别只依赖于
$A_\infty$ 代数
$(C,\mathfrak m)$ 的同伦等价类。
对
$\mathfrak M_{\mathrm{weak}}(C,\mathfrak m)$
及
$\mathfrak M_{\mathrm{weak}}(C,\mathfrak m)
^{\mathrm{bos}}$
也有相同结论。
\end{theorem}

\begin{proof}
设 $\mathfrak f\colon C\to C'$
为两个典范滤过 $A_\infty$
代数之间的同构。
取 $C^{\mathrm{odd}}$ 的基
$\mathbf v_1,\ldots,\mathbf v_m$
及 $C^{\mathrm{even}}$ 的基
$\mathbf w_1,\ldots,\mathbf w_\ell$，
以 $s_i$（$1\le i\le m$）
和 $h_j$（$1\le j\le\ell$）
为相应形式坐标。
同样地，对 $C'$ 取基
$\mathbf v'_1,\ldots,\mathbf v'_{m'}$、
$\mathbf w'_1,\ldots,\mathbf w'_{\ell'}$，
并令 $s'_i,h'_j$
为相应形式坐标。
设
$\pi'_j\colon C'\to
\Lambda_{0,\mathrm{nov}}$
（$1\le j\le m'+\ell'$）
为对后一组基的坐标投影。
\translatornote{原书实体页 343
在选定 $C'$ 的
$m'+\ell'$ 个基向量后，
仍把 $\pi'_j$ 的范围
写为 $1\le j\le m+\ell$，
并把 $\mathfrak F(h'_j)$ 的指标
写为 $j+m$；
这里均按带撇对象的基维数书写，
见疑误 A-0316。}
定义
\begin{equation}
\begin{aligned}
\mathfrak F_j
&=
(\pi'_j\circ\mathfrak f)
\Bigl(
\exp(
s_1\mathbf v_1+\cdots+s_m\mathbf v_m
+h_1\mathbf w_1+\cdots+h_\ell\mathbf w_\ell
)
\Bigr)\\
&\in
\Lambda_{0,\mathrm{nov}}
[[s_1,\ldots,s_m;h_1,\ldots,h_\ell]].
\end{aligned}
\tag{5.4.13}\label{eq:5-4-13}
\end{equation}
其意义与 (5.4.10) 相同。
令
\[
\begin{aligned}
\mathcal R
&=\Lambda_{0,\mathrm{nov}}
[[s_1,\ldots,s_m;h_1,\ldots,h_\ell]],\\
\mathcal R'
&=\Lambda_{0,\mathrm{nov}}
[[s'_1,\ldots,s'_{m'};
h'_1,\ldots,h'_{\ell'}]].
\end{aligned}
\]
定义 $\Lambda_{0,\mathrm{nov}}$-代数同态
$\mathfrak F\colon\mathcal R'\to\mathcal R$：
令
$\mathfrak F(s'_j)=\mathfrak F_j$、
$\mathfrak F(h'_j)=\mathfrak F_{j+m'}$，
再延拓为代数同态。
此外，$(C,\mathfrak m)$
的超仓西映射定义自同态
$\mathfrak P\colon\mathcal R\to\mathcal R$，
它由
$\mathfrak P(s_j)=P_j$、
$\mathfrak P(h_j)=P_{j+m}$
给出。
对 $(C',\mathfrak m')$
同样定义
$\mathfrak P'\colon\mathcal R'\to\mathcal R'$。
由于 $\mathfrak f$ 是
$A_\infty$ 同态，
下列图式交换：
\begin{center}
\[
\begin{array}{ccc}
\mathcal R'
&\xrightarrow{\ \mathfrak P'\ }&
\mathcal R'\\[2mm]
\big\downarrow{\mathfrak F}
&&\big\downarrow{\mathfrak F}\\[2mm]
\mathcal R
&\xrightarrow{\ \mathfrak P\ }&
\mathcal R .
\end{array}
\]
\renewcommand{\thefigure}{5.4.1}
\refstepcounter{figure}
{\normalfont\bfseries 图式~\thefigure}
\label{diag:5-4-1}
\end{center}
因此诱导环同态
$\mathfrak F\colon
\mathfrak M(C',\mathfrak m')
\to\mathfrak M(C,\mathfrak m)$。
显然，对应
$\mathfrak f\mapsto\mathfrak F$
是函子性的，
从而定理得证。
\end{proof}

\begin{remark}{5.4.14}\label{rem:5-4-14}
超仓西映射是形式幂级数，
所以不能径自以有理数
代替形式坐标 $s_i$。
能否找到仓西映射收敛的典范模型，
看来是一个重要问题。
这里有两种收敛问题：
一是能否把 $s_i$
代入（较小的）有理数；
二是能否把 $T$
代入（较小的）非负实数。
二者在同调镜像对称研究中
看来都具有基础性意义。
前一种收敛问题
与猜想 3.6.46 有关。
\end{remark}

下面解释
$\mathfrak M(C,\mathfrak m)^{\mathrm{bos}}$
与定界余链模空间
$\widehat{\mathcal M}(C,\mathfrak m)$
的关系。
记
\[
\operatorname{Hom}\bigl(
\mathfrak M(C,\mathfrak m)^{\mathrm{bos}},
\Lambda_{0,\mathrm{nov}}
\bigr)
\]
为从
$\mathfrak M(C,\mathfrak m)^{\mathrm{bos}}$
到 $\Lambda_{0,\mathrm{nov}}$
的所有\emph{连续}
$\Lambda_{0,\mathrm{nov}}$-代数同态所成的集合。
其拓扑按如下方式定义。
形式幂级数环
$\Lambda_{0,\mathrm{nov}}[[s_1,\ldots,s_m]]$
带两个滤过：
一个来自 $\Lambda_{0,\mathrm{nov}}$
自身的滤过，
另一个由
$s_1,\ldots,s_m$ 生成的理想给出。
分别记为 $F^\lambda$ 与 $G^n$。
在该形式幂级数环上，
以
\[
F^\lambda
\Lambda_{0,\mathrm{nov}}[[s_1,\ldots,s_m]]
\cap
G^n
\Lambda_{0,\mathrm{nov}}[[s_1,\ldots,s_m]]
\qquad
(\lambda\in\mathbb R,\ n\in\mathbb Z)
\]
为基本邻域系，定义一个拓扑。
$\mathfrak M(C,\mathfrak m)^{\mathrm{bos}}$
的拓扑由此诱导。

\begin{proposition}{5.4.15}\label{prop:5-4-15}
\[
\operatorname{Hom}\bigl(
\mathfrak M(C,\mathfrak m)^{\mathrm{bos}},
\Lambda_{0,\mathrm{nov}}
\bigr)
\cong
\widehat{\mathcal M}
(C_{\mathrm{can}},\mathfrak m_{\mathrm{can}}).
\]
这里
$(C_{\mathrm{can}},\mathfrak m_{\mathrm{can}})$
是 $(C,\mathfrak m)$ 的一个典范模型。
\end{proposition}

\begin{proof}
依定义，只需讨论典范滤过
$A_\infty$ 代数
$(C,\mathfrak m)$
及其仓西映射 $P_i^{\mathrm{bos}}$。
设
$\varphi\colon
\mathfrak M(C,\mathfrak m)^{\mathrm{bos}}
\to\Lambda_{0,\mathrm{nov}}$
为连续
$\Lambda_{0,\mathrm{nov}}$-代数同态，
令 $\varphi(s_i)=b_i$。
由于 $s_i^n\to0$，
连续性给出 $b_i^n\to0$，
故
$b_i\in\Lambda_{0,\mathrm{nov}}^+$。
于是
$\varphi\mapsto
(\varphi(s_1),\ldots,\varphi(s_m))$
给出同构
\[
\operatorname{Hom}\bigl(
\Lambda_{0,\mathrm{nov}}
[[s_1,\ldots,s_m]],
\Lambda_{0,\mathrm{nov}}
\bigr)
\longrightarrow
(\Lambda_{0,\mathrm{nov}}^+)^m.
\]
另一方面，由 (5.4.9)
与定义 5.4.11 可知，
当各
$b_i\in\Lambda_{0,\mathrm{nov}}^+$
时，下列条件
\begin{equation}
P_i^{\mathrm{bos}}(b_1,\ldots,b_m)=0,
\qquad
i=m+1,\ldots,m+\ell
\tag{5.4.16}\label{eq:5-4-16}
\end{equation}
当且仅当
\[
b=b_1\mathbf v_1+\cdots+
b_m\mathbf v_m
\in\widehat{\mathcal M}(C,\mathfrak m).
\]
命题得证。
\end{proof}

定理 M
由定理 5.4.12 的
“弱”与“形变”版本，
以及命题 5.4.15 推出
（见第 5.4.6 小节）。

\begin{remark}{5.4.17}\label{rem:5-4-17}
除了考虑取值于万有诺维科夫环的
连续同态，
还可考虑任意
Artin 环 $\mathfrak A$
上的
$\operatorname{Hom}
(\mathfrak M(C,\mathfrak m)^{\mathrm{bos}},
\mathfrak A^+)$，
其中 $\mathfrak A^+$
是 $\mathfrak A$ 的极大理想；
或者考虑完备赋值环
$\mathfrak A$，
其上给定连续同态
$\Lambda_{0,\mathrm{nov}}\to\mathfrak A$。
这样，上述讨论就更接近
Artin 环上的形变函子理论
\textbf{[Schl68]}；
另见 \textbf{[Fuk03II]}。
\end{remark}
